\فصل{کارهای پیشین و تحلیل سامانه‌های موجود}

در فصل سوم، پیشینه‌ی تحقیق و کارهای انجام‌شده در حوزه‌ی سامانه‌های پاسخ‌گویی هوشمند و جست‌وجوی مبتنی بر مدل‌های زبانی به تفصیل مورد بررسی قرار می‌گیرد. 
مبنای نظری این سامانه‌ها را می‌توان در معماری تولید افزوده با بازیابی\LTRfootnote{Retrieval-Augmented Generation (RAG)} مشاهده کرد؛ چارچوبی که دانش پارامتری مدل زبانی را با حافظه‌ی غیرپارامتری مبتنی بر اسناد خارجی ترکیب می‌کند و برای پرسش‌های دانش‌محور عملکرد بهتری نسبت به مدل‌های صرفاً مولد نشان می‌دهد~\مرجع{lewis2020rag}. در ادبیات جدید، \lr{RAG} به‌عنوان یکی از مهم‌ترین راهکارها برای کاهش نرخ توهم\LTRfootnote{Hallucination}، رفع مشکل دانش منسوخ و افزایش قابلیت ردیابی پاسخ‌های مدل‌های زبانی معرفی شده است~\مرجع{gao2024retrieval}.

هدف از این فصل، بررسی جامع و تحلیلی سامانه‌های پیشگام در این حوزه است. در این راستا، ابتدا سه مورد از اثرگذارترین سیستم‌های فعلی یعنی پرپلکسیتی\LTRfootnote{Perplexity AI}، نمای کلی هوش مصنوعی گوگل\LTRfootnote{Google AI Overviews} و مایکروسافت کوپایلت\LTRfootnote{Microsoft Bing Copilot} از منظر تاریخچه، هویت و معماری فنی تحلیل می‌شوند. 
در نهایت، با ارائه‌ی یک تحلیل تطبیقی از این سامانه‌ها، شکاف‌های پژوهشی موجود شناسایی شده و ضرورت پیاده‌سازی لایه‌ی بازرتبه‌بندی\LTRfootnote{Reranking} در سامانه‌ی پیشنهادی این پژوهش تبیین می‌گردد.

\قسمت{تحلیل سامانه‌ی پرپلکسیتی (پارادایم موتور پاسخ‌دهی)}

سامانه‌ی پرپلکسیتی در سال ۲۰۲۲ با هدف بازتعریف مفهوم جست‌وجو تأسیس شد. فلسفه‌ی بنیادین این سیستم، گذار از مفهوم سنتی «موتور جست‌وجو»\LTRfootnote{Search Engine} به «موتور پاسخ‌دهی»\LTRfootnote{Answer Engine} است. در حالی که موتورهای سنتی کاربر را به سمت منابع و پیوندها هدایت می‌کنند، پرپلکسیتی با اتخاذ رویکرد \مهم{سنتز اطلاعات}، یک پاسخ نهایی بر پایه‌ی اطلاعات بلادرنگ تولید کرده و منابع را صرفاً به‌عنوان ابزاری برای تأیید صحت\LTRfootnote{Verification} و شفافیت ارائه می‌دهد~\مرجع{aws2023perplexity}.

پیش‌زمینه‌ی این رویکرد پیش‌تر در پژوهش \lr{WebGPT} نیز مطرح شده بود که نشان داد اتصال مدل زبانی به محیط مرور وب و الزام آن به جمع‌آوری منابع در حین پاسخ‌گویی، می‌تواند شفافیت و قابلیت راستی‌آزمایی پاسخ‌ها را به شکل چشمگیری افزایش دهد~\مرجع{nakano2021webgpt}. با توجه به مستندات عمومی درباره‌ی پرپلکسیتی و نیز الگوهای رایج در سامانه‌های \lr{RAG} پیشرفته، می‌توان معماری آن را به‌صورت یک خط لوله‌ی چندمرحله‌ای به شرح زیر تحلیل کرد:

\شروع{فقرات}

\فقره
\مهم{تحلیل نیت پرسش\LTRfootnote{Query Intent Parsing}}: تحلیل معنایی پرسش توسط یک مدل زبانی بزرگ\LTRfootnote{LLM} برای تشخیص نیت کاربر.

\فقره
\مهم{بسط عاملی پرسش\LTRfootnote{Agentic Query Expansion}}: تبدیل پرسش‌های پیچیده به چندین زیرپرسش\LTRfootnote{Sub-queries} جهت حداکثرسازی پوشش اطلاعاتی.

\فقره
\مهم{بازیابی ترکیبی\LTRfootnote{Hybrid Retrieval}}: استفاده‌ی هم‌زمان از استراتژی‌های بازیابی. بازیابی تُنُک\LTRfootnote{Sparse Retrieval} (مانند \lr{BM25}) در تطبیق واژگان دقیق و نام‌ها بسیار قدرتمند است~\مرجع{robertson2009bm25}، در حالی که بازیابی چگال\LTRfootnote{Dense Retrieval} شباهت معنایی و بازنویسی‌های مفهومی را بهتر درک می‌کند~\مرجع{karpukhin2020dpr}. ترکیب این دو، نرخ بازیابی (Recall) را در مواجهه با وب به حداکثر می‌رساند.

\فقره
\مهم{بازرتبه‌بندی چندلایه\LTRfootnote{Multi-Tier Reranking}}: عبور نتایج از فیلترهایی نظیر بررسی قابلیت استخراج از \lr{HTML}، استفاده از مدل‌های رمزگذار متقاطع\LTRfootnote{Cross-Encoder} و اولویت‌دهی به اعتبار منابع.

\فقره
\مهم{تولید مستند\LTRfootnote{Grounded Generation}}: تولید پاسخ توسط مدل زبانی با الزام به درج ارجاعات به‌صورت شماره‌گذاری‌شده در متن.

\پایان{فقرات}


\قسمت{تحلیل سامانه‌ی گوگل (معماری \lr{SGE} و ادغام گراف دانش)}

سامانه‌ی نمای کلی هوش مصنوعی گوگل (\lr{AI Overviews})، به‌عنوان لایه‌ای مولد بر فراز زیرساخت جست‌وجوی این شرکت معرفی شده است که قابلیت‌های مدل‌های خانواده‌ی جِمینای\LTRfootnote{Gemini}، از جمله استدلال چندمرحله‌ای و پردازش چندوجهی را با عظیم‌ترین نمایه‌ی وب جهان ترکیب می‌کند~\مرجع{google2024io}.

یکی از برتری‌های ساختاری گوگل نسبت به سامانه‌های مستقل، استفاده از \مهم{گراف دانش}\LTRfootnote{Knowledge Graph} است. این گراف که موجودیت‌های جهان واقعی و روابط میان آن‌ها را مدل می‌کند، نقش حیاتی در برطرف کردن ابهامات معنایی\LTRfootnote{Disambiguation} و پیشنهاد مسیرهای اکتشافی دارد~\مرجع{singhal2012knowledge}.
علاوه بر این، گوگل در مستندات رسمی خود بیان می‌کند که این سامانه ممکن است از تکنیک بسطِ انشعابی\LTRfootnote{Query Fan-out} استفاده کند؛ در این تکنیک، پرسش کاربر به چند جست‌وجوی مرتبط در زیرموضوع‌ها و منابع داده‌ای مختلف گسترش می‌یابد تا پاسخ نهایی بر شواهد گسترده‌تری متکی باشد~\مرجع{google2024searchcentral}. 

برای ادغام نتایج حاصل از این بازیابی‌های موازی، یکی از روش‌های شناخته‌شده در ادبیات بازیابی اطلاعات، الگوریتم \مهم{ادغام رتبه‌ی متقابل}\LTRfootnote{Reciprocal Rank Fusion (RRF)} است~\مرجع{cormack2009reciprocal}. امتیاز هر سند در این الگوریتم به صورت زیر محاسبه می‌شود:

\شروع{equation}
\mathrm{Score}(d) = \sum_{q \in \mathrm{Queries}} \frac{1}{k + \mathrm{rank}(q, d)}
\پایان{equation}

در این رابطه‌، $k$ یک مقدار ثابت (معمولاً برابر با ۶۰) است. این روش به هر سند بر اساس رتبه‌ی آن امتیازی معکوس اختصاص داده و اسنادی را که در جست‌وجوهای مختلف رتبه‌ی بالایی کسب کرده‌اند، تقویت می‌کند.


\قسمت{تحلیل مایکروسافت کوپایلت (معماری پرومتئوس)}

سامانه‌ی مایکروسافت کوپایلت، نتیجه‌ی ادغام نمایه‌ی موتور جست‌وجوی بینگ با مدل‌های خانواده‌ی \lr{GPT} شرکت \lr{OpenAI} است. مایکروسافت لایه‌ای اختصاصی به نام \مهم{پرومتئوس}\LTRfootnote{Prometheus} را به‌عنوان میانجی توسعه داده است. در این معماری، بخشی به نام هماهنگ‌کننده\LTRfootnote{Bing Orchestrator} پرسش‌های داخلی مرتبط را به‌صورت تکراری تولید کرده و فرآیندی به نام اتصال\LTRfootnote{Grounding} را شکل می‌دهد که طی آن مدل بر داده‌های به‌روز بینگ استدلال می‌کند~\مرجع{ribas2023bing}.

از منظر پژوهشی، این معماری را می‌توان مبتنی بر حلقه‌ی استدلال-بازیابی دانست که با کار پژوهشگرانی نظیر تریودی و همکاران (\lr{IRCoT}) نیز پشتیبانی می‌شود. این الگو نشان می‌دهد که در پرسش‌های پیچیده و چندمرحله‌ای، بازیابی یک‌باره کافی نیست؛ بلکه مدل باید پس از هر گامِ استدلال، بر اساس نیاز اطلاعاتی جدید جست‌وجوی تازه‌ای انجام دهد تا نقص‌های اطلاعاتی برطرف گردد~\مرجع{trivedi2023interleaving}.


\قسمت{تحلیل تطبیقی سامانه‌های پیشرو}

برای درک بهتر تفاوت‌ها و شباهت‌های بنیادین میان سامانه‌های بررسی‌شده، ویژگی‌های فنی و معماری احتمالی آن‌ها بر اساس مستندات عمومی در جدول~\رجوع{جدول:تحلیل-تطبیقی} خلاصه‌سازی شده است. این ماتریس نشان می‌دهد که چگونه تفاوت در فلسفه‌ی طراحی، به استراتژی‌های متفاوت در توسعه‌ی \lr{RAG} تجاری منجر شده است.

\شروع{لوح}[t]
\وسط‌چین
\small
\شرح{ماتریس فنی و تحلیل تطبیقی معماری سامانه‌های پاسخ‌گویی هوشمند}
\شروع{جدول}{|p{3cm}|p{3.5cm}|p{3.5cm}|p{3.5cm}|}
\خط‌پر
ویژگی فنی & پرپلکسیتی & گوگل & مایکروسافت \\
\خط‌پر
\خط‌پر
فلسفه طراحی & موتور پاسخ‌دهی پژوهش‌محور & جست‌وجوی مولد مبتنی بر نمایه‌ی \lr{Search} & دستیار مبتنی بر جست‌وجو و بهره‌وری \\
شواهد عمومی معماری & پاسخ سنتزشده همراه با ارجاع دقیق & ترکیب \lr{Gemini} و سیستم‌های سنتی جست‌وجو & معماری پرومتئوس + \lr{Orchestrator} \\
استراتژی احتمالی بازیابی & \lr{RAG} وب‌محور (احتمالاً ترکیبی) & تکنیک \lr{Fan-out} (تأییدشده در مستندات) & تولید پرسش‌های داخلی تکرارشونده \\
ادغام و کنترل شواهد & قابل تحلیل با الگوی بازرتبه‌بندی & ادغام نتایج انشعابی (الگوهای مشابه \lr{RRF}) & اتصال متوالی روی داده‌های بینگ \\
نحوه‌ی ارجاع‌دهی & ارجاعات دقیق و شماره‌گذاری‌شده & لینک‌های پشتیبان متنی (کارت‌های منبع) & ارجاعات درون‌متنی تعاملی \\
\خط‌پر
\پایان{جدول}
\برچسب{جدول:تحلیل-تطبیقی}
\پایان{لوح}


\قسمت{تحلیل تئوریک بازرتبه‌بندی و تبیین شکاف پژوهشی}

در سامانه‌های \lr{RAG}، کیفیت پاسخ نهایی تنها تابع توانمندی مدل زبانی نیست، بلکه به‌شدت به کیفیت شواهد بازیابی‌شده وابسته است. اگر مرحله‌ی بازیابی، اسناد نامرتبط یا متناقض را وارد بافتار کند، مدل زبانی ممکن است پاسخی روان اما نادرست تولید کند. برای کنترل این چالش، فرآیند بازیابی اطلاعات معمولاً در دو لایه انجام می‌شود:

\شروع{فقرات}

\فقره
\مهم{مرحله‌ی اول (رمزگذارهای دوگانه)\LTRfootnote{Bi-Encoders}}: پرسش و اسناد مستقل از هم به بردار تبدیل شده و شباهت آن‌ها ارزیابی می‌شود. مزیت اصلی این لایه سرعت بسیار بالا است (مانند معماری \lr{Sentence-BERT} که زمان جست‌وجو را به شدت کاهش می‌دهد~\مرجع{reimers2019sentence})، اما دقت پایینی در تشخیص روابط پیچیده‌ی معنایی دارد.

\فقره
\مهم{مرحله‌ی دوم (رمزگذارهای متقاطع)\LTRfootnote{Cross-Encoders}}: در این لایه که به‌عنوان بازرتبه‌بندی شناخته می‌شود، پرسش و سند هم‌زمان وارد شبکه‌ی ترانسفورمر می‌شوند تا تعاملات معنایی آن‌ها توکن‌به‌توکن محاسبه گردد. پژوهش‌ها نشان می‌دهد استفاده از مدل‌هایی نظیر \lr{BERT} به عنوان بازرتبه‌بند می‌تواند معیارهای رتبه‌بندی را به شکل قابل‌توجهی ارتقا دهد~\مرجع{nogueira2019passage}. همچنین مدل‌های تعاملِ دیرهنگام\LTRfootnote{Late Interaction} نظیر \lr{ColBERT} نیز به‌عنوان راهکار میانی برای ایجاد تعادل بین سرعت و دقت پیشنهاد شده‌اند~\مرجع{khattab2020colbert}.

\پایان{فقرات}

\زیرقسمت{RAG تجاری در برابر RAG آکادمیک و شکاف‌های موجود}
سامانه‌های تجاریِ بررسی‌شده در این فصل، نمونه‌هایی از \lr{RAG} عملیاتی در مقیاس وب هستند. تفاوت اساسی آن‌ها با نمونه‌های دانشگاهی در مواجهه با محیط پویای وب نهفته است. در پژوهش‌های دانشگاهی مجموعه‌داده‌ها غالباً ایستا و تمیز هستند، اما در وب واقعی، سیستم با محتوای تغییرپذیر، اطلاعات متناقض و نویز شدید روبه‌رو است که نقش بازرتبه‌بندی را حیاتی می‌سازد.

با وجود این اهمیت، بررسی ادبیات موضوع حاکی از چند شکاف پژوهشی است: نخست آنکه سامانه‌های تجاری جزئیات کمّی و قابل‌بازتولیدی از تأثیر لایه‌ی بازرتبه‌بندی منتشر نمی‌کنند. 
دوم آنکه پژوهش‌هایی نظیر \lr{FreshLLMs} اثبات کرده‌اند که مدل‌های زبانی در مواجهه با دانشی که به‌سرعت تغییر می‌کند دچار ضعف جدی می‌شوند و تزریق شواهد کاملاً به‌روز از وب ضروری است~\مرجع{vu2023freshllms}. 
از سوی دیگر، پژوهش \lr{Lost in the Middle} نشان می‌دهد که نحوه‌ی چیدمان و موقعیت قرارگیری شواهد در بافتار ارائه‌شده به مدل، مستقیماً بر کیفیت خروجی اثرگذار است و مدل‌ها اطلاعاتِ میانه‌ی بافتار را نادیده می‌گیرند؛ موضوعی که ضرورت بازرتبه‌بندی دقیق را دوچندان می‌کند~\مرجع{liu2024lost}.

بنابراین، توجیه علمی و هدف اصلی پژوهش حاضر آن است که با طراحی یک سیستم کنترل‌شده و اجرای مطالعه‌ی حذف\LTRfootnote{Ablation Study} بر روی داده‌های پویای وب، این شکاف پژوهشی را پوشش دهد. در این راستا، برای ارزیابی دقیق‌تر علاوه بر معیارهای سنتی بازیابی، می‌توان از چارچوب‌های نوین ارزیابی نظیر \lr{RAGAS} بهره برد تا مشخص شود پاسخ نهایی تا چه حد به شواهد ارائه‌شده وفادار (\lr{Faithful}) بوده است~\مرجع{es2023ragas}. این تحلیل تجربی به صورت مستدل اثبات خواهد کرد که لایه‌ی بازرتبه‌بندی تا چه میزان قادر است نرخ توهم مدل را در محیط واقعی کاهش دهد.